\documentclass[11pt]{article}

% ---------- Packages ----------
\usepackage[a4paper,margin=1in]{geometry}
\usepackage{subfiles}
\usepackage{amsmath,amssymb,bm}
\usepackage[dvipsnames]{xcolor}
\usepackage{hyperref}
\usepackage[numbers,sort&compress]{natbib}
\usepackage{booktabs}
\usepackage{microtype}

\hypersetup{
  colorlinks=true,
  linkcolor=blue!60!black,
  citecolor=blue!60!black,
  urlcolor=blue!60!black
}

% ---------- Shorthand ----------
\newcommand{\Mpl}{M_{\rm Pl}}
\newcommand{\aB}{\alpha_{\rm B}}
\newcommand{\aK}{\alpha_{\rm K}}
\newcommand{\aM}{\alpha_{\rm M}}
\newcommand{\aT}{\alpha_{\rm T}}
\newcommand{\azero}{a_0}
\newcommand{\LCDM}{$\Lambda$CDM}

\title{Appendix 9 (to main theory): CMB and Linear-Cosmology Consistency}
\author{}
\date{}
\begin{document}
\let\phi\varphi
\maketitle

\begin{abstract}
We show that the bi-conformal scalar--tensor framework defined in the main paper \cite{Lotishvili2026} satisfies current linear cosmological constraints.
On a spatially flat Friedmann--Lema\^{\i}tre--Robertson--Walker (FLRW) background
written in the physical cosmic-time frame tied to matter clocks, the homogeneous
scalar mode satisfies $\phi_0=0$ as a consistency requirement of simultaneously
(i) using comoving proper time $t=\tau$ and (ii) imposing the bi-conformal ansatz
for $g_{tt}$. Under this condition all four Bellini--Sawicki EFT functions vanish
identically ($\aK=\aB=\aM=\aT=0$), so the linear Einstein--Boltzmann system reduces
to that of \LCDM{} with the same matter content and initial conditions. Within this
metric-sector limit, the primary CMB angular power spectrum, BAO scale, supernova
distance measures, and BBN expansion-rate constraints are inherited at linear order;
possible matter-sector pressure-history backreaction is a separate issue outside
the linear claim made here. The distinguishing
prediction of the framework lies in the quasi-static, weak-field regime around
bound structures: the MOND sector implies a redshift evolution
$\azero(z)=cH(z)/(2\pi)$.
\end{abstract}

% ============================================================
\section{Purpose and scope}
\label{sec:scope}
% ============================================================

This appendix is restricted to the covariant action of Ref.~\cite{Lotishvili2026}
and its FLRW / linear-perturbation consequences. It is \emph{not} an exploration
of alternative redshift interpretations in static backgrounds.

The objective is to demonstrate that (i) the background expansion is standard
and (ii) linear perturbations obey the same Einstein--Boltzmann evolution as in
\LCDM{}, while (iii) the new phenomenology (MOND sector) appears only in the
nonlinear, quasi-static limit around gravitationally bound systems.

% ============================================================
\section{Framework and conventions}
\label{sec:conventions}
% ============================================================

We adopt:
\begin{itemize}
\item Metric signature $(-,+,+,+)$.
\item Horndeski kinetic variable
\begin{equation}
X \equiv -\frac{1}{2}\,g^{\mu\nu}\partial_\mu\phi\,\partial_\nu\phi\,.
\end{equation}
\item Single-field minimally coupled Horndeski subset with constant $G_4=\Mpl^2/2$,
$G_3=G_5=0$, and a kinetic sector fixed by the bi-conformal consistency of the field
equations (see Secs.~II.B and VI of Ref.~\cite{Lotishvili2026} for the derivation and
stability analysis).
\item Minimal matter coupling: the matter action depends on $(g_{\mu\nu},\psi)$ only
(no direct $\phi$--matter coupling and no disformal matter coupling).
\item EFT-of-DE mapping and $\alpha$-functions follow Bellini \& Sawicki \cite{BelliniSawicki2014}.
\end{itemize}

\paragraph{Note on the kinetic-term sign.}
For the cosmological results in this appendix the sign choice is immaterial because
$\phi_0=0$ implies that the scalar contributes neither background energy density nor
pressure, and it does not source metric potentials at linear order. Full stability
conditions for the kinetic sector are provided in Ref.~\cite{Lotishvili2026}.

% ============================================================
\section{Action and field equations}
\label{sec:action_app9}
% ============================================================

The gravitational sector is defined by the action (Eq.~(3) of Ref.~\cite{Lotishvili2026})
\begin{equation}
S = \frac{1}{16\pi G}\int d^4x\,\sqrt{-g}\,
\Bigl[R + \tfrac{1}{2}\,g^{\mu\nu}\partial_\mu\phi\,\partial_\nu\phi\Bigr]
+ S_{\rm m}[g_{\mu\nu},\psi]\,.
\label{eq:action_app9}
\end{equation}

Direct variation of \eqref{eq:action_app9} with respect to $g^{\mu\nu}$ yields the metric equation
\begin{equation}
G_{\mu\nu} = 8\pi G\,T^{(\rm m)}_{\mu\nu}-\tfrac{1}{2}\,\tau^{(\phi)}_{\mu\nu},
\label{eq:einstein_app9}
\end{equation}
where the canonical scalar stress--energy tensor is
\begin{equation}
\tau^{(\phi)}_{\mu\nu}
= \partial_\mu\phi\,\partial_\nu\phi
-\tfrac{1}{2}g_{\mu\nu}\,g^{\alpha\beta}\partial_\alpha\phi\,\partial_\beta\phi\,.
\label{eq:Tphi_app9}
\end{equation}
Direct $\phi$-variation under minimal coupling ($\delta S_{\rm m}/\delta\phi\equiv 0$) gives the vacuum scalar equation $\square\phi=0$ also in the presence of matter. The reduced-sector trace consistency of \eqref{eq:einstein_app9} after the bi-conformal ansatz (Appendix~1, §Scalar field equation; Appendix~3) gives the source relation used throughout the main theory,
\begin{equation}
\square\phi = -\frac{8\pi G}{c^4}\,T\,,
\qquad T\equiv g^{\mu\nu}T^{(\rm m)}_{\mu\nu},
\label{eq:eoms}
\end{equation}
which is not an independent Euler--Lagrange equation.

% ============================================================
\section{FLRW background and the $\phi_0=0$ consistency in cosmic time}
\label{sec:flrw_app9}
% ============================================================

Consider a spatially flat FLRW background written in cosmic time,
\begin{equation}
ds^2 = -c^2\,dt^2 + a^2(t)\,\delta_{ij}\,dx^i dx^j\,,
\label{eq:flrw}
\end{equation}
where $t=\tau$ is the proper time measured by matter-comoving observers.

In the same coordinate frame the bi-conformal ansatz of Ref.~\cite{Lotishvili2026}
imposes $g_{tt}=-e^{\phi}$. Requiring simultaneous satisfaction of
(i) the physical cosmic-time condition $g_{tt}=-1$ in Eq.~\eqref{eq:flrw} and
(ii) the bi-conformal form therefore gives
\begin{equation}
e^{\phi_0}=1 \quad\Rightarrow\quad \phi_0=0\,,
\label{eq:phi0}
\end{equation}
on the homogeneous FLRW background.

\paragraph{Escape hatch (for clarity).}
One may work in a non-comoving time parametrization with $t\neq\tau$ and keep
$\phi_0\neq0$, but then the background is not written in the standard FLRW
cosmic-time gauge tied to matter clocks. The present appendix formulates the
consistency statement in the physical cosmic-time frame.


\paragraph{Clarifying remark (construction vs dynamics).}
Equation~\eqref{eq:phi0} shows that, in the matter-clock cosmic-time frame, the homogeneous
cosmological metric scalar is \emph{fixed} to vanish by simultaneous imposition of the FLRW
cosmic-time condition and the bi-conformal $g_{tt}$ ansatz. Accordingly, the agreement with GR at
the background level is a \emph{consistency consequence of the construction}, rather than an
independent dynamical prediction. This statement concerns the homogeneous metric mode
$\phi_0$, not the coarse-grained pressure-history diagnostic used in the process-time sector.

\paragraph{Immediate consequence.}
When $\dot\phi_0=0$ the scalar stress tensor vanishes on the background:
\begin{equation}
\dot\phi_0=0 \ \Rightarrow\  \partial_\mu\phi_0=0 \ \Rightarrow\ 
\tau^{(\phi)}_{\mu\nu}=0\,,
\end{equation}
so the background expansion follows the standard Friedmann equations for the
same matter content as in \LCDM{}. Explicitly,
\begin{equation}
3\Mpl^2 H^2 = \rho,\qquad 2\Mpl^2\dot H = -(\rho+p),\qquad 2\Mpl^2\dot H + 3\Mpl^2 H^2 = -p,
\label{eq:friedmann_app9}
\end{equation}
where $H\equiv \dot a/a$ and $(\rho,p)$ denote the total matter--radiation energy
density and pressure.

% ============================================================
\section{EFT mapping: vanishing $\alpha$-functions}
\label{sec:eft}
% ============================================================

In the Bellini--Sawicki EFT-of-DE parametrization \cite{BelliniSawicki2014}, linear
modifications to scalar and tensor perturbations enter through the four functions
$\{\aK,\aB,\aM,\aT\}$. For the present model the Horndeski functions satisfy
$G_4=\Mpl^2/2=\mathrm{const}$, $G_3=G_5=0$, and the homogeneous mode obeys
$\dot\phi_0=0$ in the physical cosmic-time frame (Sec.~\ref{sec:flrw_app9}). Using the
standard definitions of Ref.~\cite{BelliniSawicki2014} (see also the detailed
derivation in Ref.~\cite{Lotishvili2026}), one obtains the following minimal
substitution proof:
\begin{itemize}
\item \textbf{Tensor speed}: $\aT$ receives contributions from $G_{4,X}$ and $G_5$. Here
$G_{4,X}=0$ and $G_5=0$, hence $\aT=0$ (consistent with the GW+GRB speed bound
from GW170817/GRB~170817A~\cite{LIGO_GW170817}).
\item \textbf{Planck-mass running}: $\aM\equiv \frac{d\ln M_*^2}{d\ln a}$ with
$M_*^2=2G_4$. Since $G_4=\Mpl^2/2=\mathrm{const}$, $M_*^2=\Mpl^2$ and $\aM=0$.
\item \textbf{Braiding}: $\aB$ is sourced by $G_3$ and by $\phi$-dependence of $G_4$.
Here $G_3=0$ and $G_{4,\phi}=0$, hence $\aB=0$.
\item \textbf{Kineticity}: $\aK$ is proportional to the background kinetic energy
$X\propto \dot\phi_0^{\,2}$, with sign fixed by the Horndeski map $G_2=-(\Mpl^2/2)X$ (consistent with the main-text and Appendix~11 normalization; the shorthand $G_2=-X/2$ used in some places of this appendix sets $\Mpl=1$).
Since $\dot\phi_0=0$ (Sec.~\ref{sec:flrw_app9}), $X=0$ and $\aK=0$ trivially.
For $\dot\phi_0\neq 0$, $\aK=-X/H^2<0$, exposing the phantom kinetic sign and
requiring the EFT completion of Appendix~11 for stability.
\end{itemize}
Consequently,
\begin{equation}
\aK=\aB=\aM=\aT=0\,.
\label{eq:alphazero}
\end{equation}
Hence, under the construction choice $\dot\phi_0=0$ (Sec.~\ref{sec:flrw_app9}),
the linear gravitational sector is GR-identical at leading order in the metric
potentials.

% ============================================================
\section{Linear perturbations: Poisson equation and slip}
\label{sec:linear}
% ============================================================

In Newtonian gauge,
\begin{equation}
ds^2 = -(1+2\Psi)c^2dt^2 + a^2(t)(1-2\Phi)\delta_{ij}dx^i dx^j\,,
\end{equation}
the GR-identical EFT limit \eqref{eq:alphazero} implies the standard relations
\begin{equation}
\Phi-\Psi = 0,
\qquad
-k^2\Psi = 4\pi G a^2 \rho\,\Delta,
\label{eq:poisson_app9}
\end{equation}
with $\Delta$ the comoving density contrast.
In EFT language, departures from these GR forms arise through combinations of
$\aB,\aM,\aT$; since they vanish here, the linear Einstein--Boltzmann system for
photons, baryons, CDM and neutrinos is unchanged (for the same matter content
and initial conditions). The scalar perturbation propagates as a matter-sourced
mode.  Since $\partial_\mu\phi_0=0$ on the background, the scalar stress tensor
has no term linear in $\delta\phi$; in this reduced metric sector the sourced
scalar therefore does not back-react on the metric potentials at linear order.

\paragraph{Matter-sourced scalar mode.}
In Fourier space the linear scalar perturbation obeys a damped wave equation on
FLRW, sourced by the matter trace (Appendix~1 sourced EOM):
\begin{equation}
\delta\ddot\phi_k + 3H\,\delta\dot\phi_k + \frac{c^2k^2}{a^2}\,\delta\phi_k
= \frac{8\pi G}{c^2}\,\delta T_k,\qquad \delta T\equiv \delta\bigl(g^{\mu\nu}T^{(\rm m)}_{\mu\nu}\bigr),
\label{eq:deltaphi_app9}
\end{equation}
illustrating that the scalar mode is sourced by $\delta T_k$ while remaining
decoupled from the metric potentials at linear order under the
$\partial_\mu\phi_0=0$ and $\alpha_i=0$ conditions. The
term ``spectator'' is used here in the restricted sense of metric-decoupling,
not full dynamical decoupling: the mode is driven by matter and its kinetic
stability inherits the phantom-sign caveat of Appendix~11.

% ============================================================
\section{Observational consequences and regime separation}
\label{sec:obs}
% ============================================================

\paragraph{Inherited linear observables.}
Because both the background expansion and the linear gravitational sector are
GR-identical in the same-matter-content metric-sector limit, the following are
inherited at linear order: primary CMB anisotropies, BAO standard ruler,
supernova distance measures, and the BBN expansion-rate constraint $H(T)$.
This inheritance statement does not by itself quantify any matter-sector
backreaction associated with the pressure-history diagnostic
$\phi_{\mathrm{bg}}(z)$.


\paragraph{From a GR-identical FLRW background to MOND in galaxies (sketch).}
With $\phi_0=0$ the scalar is absent in the homogeneous cosmological background, but it can be
\emph{inhomogeneously} sourced by matter through $\delta T$ as in Eq.~\eqref{eq:deltaphi_app9}. The MOND
phenomenology is claimed to arise only after taking the weak-field, quasi-static limit around a
bound structure (sub-horizon $k\gg aH/c$) and solving the \emph{nonlinear} field equations sourced by
localized overdensities. In this limit time derivatives are parametrically suppressed relative to
spatial gradients, schematically $\partial_t \sim H \ll c\,\nabla$, and the scalar sector contributes
an additional, environment-dependent term to the effective gravitational law. A fully rigorous
derivation of the MOND limit from the nonlinear equations is presented in the main paper
(Ref.~\cite{Lotishvili2026}); the purpose of the present appendix is only to document that the same
construction leaves the FLRW/linear Einstein--Boltzmann system unchanged.

\paragraph{Where the new physics enters.}
Deviations arise only in the nonlinear, quasi-static, weak-field regime around bound
structures, where the same field equations yield the MOND phenomenology discussed in
Ref.~\cite{Lotishvili2026}. The central falsifiable cosmological linkage is introduced by (i) defining a Hubble wavenumber and its associated wavelength via the standard Fourier relation, and (ii) \emph{identifying} the MOND acceleration scale with the inverse of that wavelength in units of $c^2$:
\begin{equation}
k_H \equiv \frac{aH}{c},\qquad
\lambda_H \equiv \frac{2\pi a}{k_H}=\frac{2\pi c}{H},\qquad
\azero \equiv \frac{c^2}{\lambda_H}\,.
\label{eq:2pi_chain}
\end{equation}
\noindent The last relation in Eq.~\eqref{eq:2pi_chain} is a \emph{physical identification} of the characteristic acceleration scale with a cosmological wavelength, not a mathematical identity.
Therefore,
\begin{equation}
\azero(z)=\frac{cH(z)}{2\pi},
\label{eq:a0z_app9}
\end{equation}
which implies a redshift drift of BTFR normalisation (in principle testable, subject to
astrophysical systematics and calibration).

Since $v_{\rm flat}^4\propto \azero GM$ in the deep-MOND limit, consistent with the
observed radial-acceleration relation \cite{McGaugh2016}, one has
\begin{equation}
v_{\rm flat}\propto \azero^{1/4},
\label{eq:v_a0_quarter}
\end{equation}
so, for example, in a baryons-only cosmology ($\Omega_b\approx0.05$) natural to this framework, $H(1)/H_0\approx 1.16$, giving a $\sim$16\% increase in $\azero$ at $z\simeq 1$ and a corresponding $\simeq 4\%$ upward shift in the BTFR velocity scale at fixed baryonic mass.

\paragraph{Possible deviations (targets for tests).}
\begin{itemize}
\item BTFR normalisation drift from Eq.~\eqref{eq:a0z_app9}.
\item Galaxy-scale weak-lensing shear profiles in the quasi-static regime.
\item Dwarf/satellite dynamics sensitivity to environment (external-field effects) as a
structural consequence of the nonlinear field equations.
\end{itemize}

% ============================================================
\section{Comparison summary}
\label{sec:tables}
% ============================================================

\begin{table}[ht]
\centering
\caption{Linear-cosmology sector: status of key observables.}
\begin{tabular}{@{}ll@{}}
\toprule
Observable & Status in ISPG bi-conformal theory \\ \midrule
Primary CMB $C_\ell$ & Unchanged in the metric-sector limit \\
BAO scale & Unchanged for the same matter content \\
SNe distances & Unchanged for the same background distances \\
BBN $H(T)$ & Unchanged in the metric-sector limit \\
Linear growth $f\sigma_8$ & Unchanged in the metric-sector limit \\
Lensing potential $(\Phi+\Psi)/2$ & Unchanged at linear order in this limit \\
\bottomrule
\end{tabular}
\end{table}

\begin{table}[ht]
\centering
\caption{Illustrative comparison to common modified-gravity frameworks (linear sector).
TeVeS references: Bekenstein \cite{Bekenstein2004}; known CMB-related tensions without
additional ingredients: see, e.g., Skordis \& Z\l{}o\'snik \cite{SkordisZlosnik2021}.}
\begin{tabular}{@{}lcccc@{}}
\toprule
Model & $\aT$ & $\aM$ & $\aB$ & CMB linear sector \\ \midrule
ISPG bi-conformal (this work) & $0$ & $0$ & $0$ & GR-identical \\
Brans--Dicke (generic) & $0$ & $\neq 0$ & $\neq 0$ & modified growth/ISW \\
$f(R)$ (generic) & $0$ & $\neq 0$ & $\neq 0$ & modified growth/ISW \\
Galileon / braiding models & $\approx 0$ & $\neq 0$ & $\neq 0$ & modified potentials \\
TeVeS-type theories & varies & varies & varies & known tensions without extra ingredients \\
\bottomrule
\end{tabular}
\end{table}

\begin{table}[ht]
\centering
\caption{Nonlinear/galaxy regime: distinguishing predictions and tests.}
\begin{tabular}{@{}ll@{}}
\toprule
Item & Status \\ \midrule
MOND phenomenology & Emergent in quasi-static limit \\
BTFR $v^4\propto \azero GM$ & Predicted; normalisation evolves via $\azero(z)$ \\
$\azero(z)=cH(z)/(2\pi)$ & Distinct cosmological linkage (testable in principle) \\
Dwarf/satellite dynamics & Environment sensitivity (EFE-type behaviour) \\
\bottomrule
\end{tabular}
\end{table}
\end{document}
